% 贯穿四点图 V={1,2,3,4}，E={12,13,23,24,34}；从1出发，邻居按编号递增访问。
% 队列轨迹：[1] -> [2,3] -> [3,4] -> [4] -> []；d=(0,1,1,2)，p(2)=p(3)=1，p(4)=2。
\begin{tikzpicture}[font=\figurefont\fontsize{8}{10}\selectfont,>=Stealth,
  vtx/.style={circle,draw=FigureStroke,line width=.8pt,minimum size=6.5mm,inner sep=0pt,text=FigureInk},
  tree/.style={draw=FigureBlue,line width=1.3pt},
  other/.style={draw=FigureMuted,line width=.7pt,dashed},
  qcell/.style={draw=FigureStroke,line width=.6pt,minimum width=5.2mm,minimum height=5.2mm,inner sep=0pt,fill=FigureBlueFill},
  note/.style={font=\FigureNoteFont\fontsize{7.5}{9}\selectfont,text=FigureMuted}]
% ---------- (a) 层与父指针 ----------
\begin{scope}
  \fill[FigureBlueFill] (-0.45,-1.55) rectangle (0.45,1.35);
  \fill[FigureBlueFill] (0.85,-1.55) rectangle (1.75,1.35);
  \fill[FigureBlueFill] (2.15,-1.55) rectangle (3.05,1.35);
  \node[note] at (0,1.6) {$d=0$};
  \node[note] at (1.3,1.6) {$d=1$};
  \node[note] at (2.6,1.6) {$d=2$};
  \node[vtx,fill=white] (v1) at (0,0) {1};
  \node[vtx,fill=white] (v2) at (1.3,0.8) {2};
  \node[vtx,fill=white] (v3) at (1.3,-0.8) {3};
  \node[vtx,fill=white] (v4) at (2.6,0) {4};
  \draw[other] (v2)--(v3) (v3)--(v4);
  \draw[tree,->] (v1)--(v2);
  \draw[tree,->] (v1)--(v3);
  \draw[tree,->] (v2)--(v4);
  \node at (1.3,-2.0) {(a) 距离层与父边（蓝色箭头）};
\end{scope}
% ---------- (b) 队列轨迹 ----------
\begin{scope}[xshift=4.3cm]
  \node[anchor=west] at (0,1.6) {出队};
  \node[anchor=west] at (1.0,1.6) {新发现};
  \node[anchor=west] at (2.25,1.6) {出队后的队列};
  \draw[FigureStroke,line width=.5pt] (0,1.35)--(4.6,1.35);
  \foreach \row/\out/\new/\q in {0/{初始}/{1}/{1},1/{1}/{2, 3}/{2,3},2/{2}/{4}/{3,4},3/{3}/{无}/{4},4/{4}/{无}/{}}{
    \pgfmathsetmacro{\y}{1.0-0.55*\row}
    \node[anchor=west] at (0,\y) {\out};
    \node[anchor=west] at (1.0,\y) {\new};
  }
  % 队列格子
  \node[qcell,fill=FigureGreenFill] at (2.5,1.0) {1};
  \node[qcell] at (2.5,0.45) {2}; \node[qcell,fill=FigureGreenFill] at (3.02,0.45) {3};
  \node[qcell,fill=FigureGreenFill] at (2.5,0.45) {2};
  \node[qcell] at (2.5,-0.1) {3}; \node[qcell,fill=FigureGreenFill] at (3.02,-0.1) {4};
  \node[qcell] at (2.5,-0.65) {4};
  \node[anchor=west,note] at (2.25,-1.2) {空，结束};
  \node[note,anchor=west] at (3.45,0.45) {$\leftarrow$入队时即标记};
  \node at (2.3,-2.0) {(b) 队首在左，绿色为本步新入队};
\end{scope}
\end{tikzpicture}
